\section{ Motivation }

\begin{frame}
	\frametitle{ Womit beschäftige ich mich? }
    \begin{itemize}
        \item Supraleitung beschrieben über BCS-Theorie 
        \item Wende auf Bienenwabengitter, Gitterstrukur z.B. von Graphen an
        \item Ziel: Das Auftreten von Supraleitung und dessen Begrenzungen in dem Bienenwabengitter zu untersuchen
    \end{itemize}
\end{frame}

% \begin{frame}
% 	\frametitle{ Argumentative derivation }
%     \begin{itemize}
%         \item coordination number $z$
%         \item Frobenius norm of mean-field commutators are suppressed $\rightarrow$ they behave classically
%         \begin{align*}
%             \frac{\norm{\left[\Va,\Vb\right]}_{\text{F}}}{\norm{\Va\Vb}_{\text{F}}} \propto \frac1{z}
%         \end{align*}
%         \item spins are only weakly correlated
%         \begin{align*}
%             \expval{\Sai(t) \Sbj(0)} &\propto \frac1{z^{\taxicab{i}{j}}}, & \taxicab{i}{j}&: \text{taxicab distance} 
%         \end{align*}
%         \item local spin environments: large number of independent random variables
%         $\rightarrow$ central limit theorem 
%         $\rightarrow$ dynamics Gaussian mean-fields
%     \end{itemize}
% \end{frame}